\documentclass[11pt,a4paper]{article}
\usepackage[utf8]{inputenc}
\usepackage[T1]{fontenc}
\usepackage[brazil]{babel}
\usepackage{lmodern}
\usepackage{amsmath,amssymb,amsthm}
\usepackage{booktabs}
\usepackage{float}
\usepackage[margin=2.5cm]{geometry}
\usepackage{pgfplots}
\pgfplotsset{compat=1.16}
\usepackage[hidelinks]{hyperref}

\newtheorem{teorema}{Teorema}
\newtheorem{proposicao}[teorema]{Proposição}
\newtheorem{lema}[teorema]{Lema}
\theoremstyle{remark}
\newtheorem{observacao}[teorema]{Observação}

\title{Uma aproximação para o perímetro da elipse\\ via harmônicos toroidais e funções de Bessel}
\author{}
\date{Novembro de 2026}

\begin{document}
\maketitle

\begin{abstract}
O perímetro da elipse é, a menos de um fator, uma função de Legendre de grau
$\tfrac12$ (harmônico toroidal) calculada em $\cosh\xi$, com $\xi=\ln(a/b)$. Partindo
dessa identidade conhecida, aplicamos a ideia das aproximações uniformes de
Olver para funções de Legendre, que as escrevem em termos das funções de Bessel
modificadas $I_0$ e $I_1$. A fórmula de dois termos resultante não tem nenhuma constante ajustada e
erra no máximo $0{,}7\%$. Uma pequena função de correção par em $\xi$,
determinada por um aproximante de Padé da sua série exata e pelo seu limite
assintótico exato, leva o erro relativo a no máximo $4{,}7\times10^{-14}$ para
$b/a\ge0{,}2$ (Ramanujan II: $8{,}4\times10^{-7}$) e a $10^{-21}$ em $b/a=0{,}5$.
A fórmula supera Ramanujan II para todo $b/a\gtrsim10^{-6}$.
\end{abstract}

\section{A fórmula}

Sejam $a>b>0$ os semieixos e
\begin{equation}
\xi=\ln\frac ab,\qquad t=\xi^2,\qquad
B(\xi)=\frac{\xi\coth\xi-1}{8\xi^2}.
\end{equation}
A aproximação proposta é
\begin{equation}\label{eq:formula}
\boxed{\;
P\approx2\pi\sqrt{ab}\,\sqrt{\frac{\xi}{\sinh\xi}}\;
\Big[\,I_0(\xi)+\xi\,I_1(\xi)\,B(\xi)\,\big(1+\varphi(t)\big)\Big],
\qquad
\varphi(t)=\frac{\sum_{k=1}^{5}p_kt^k}{1+\sum_{k=1}^{5}q_kt^k}\;}
\end{equation}
em que $I_0$ e $I_1$ são as funções de Bessel modificadas de primeira espécie. As
constantes estão na Tabela~\ref{tab:c}. No círculo ($a=b$), usa-se o limite $P=2\pi a$.

\begin{table}[H]
\centering
\begin{tabular}{lr@{\qquad}lr}
\toprule
$p_1$ & $7/320=0.021875$ & $q_1$ & $0.25525899377591932$\\
$p_2$ & $0.0027421486237688702$ & $q_2$ & $0.024338672837890872$\\
$p_3$ & $8.0863833392703183\times10^{-5}$ & $q_3$ & $0.0010607656517810025$\\
$p_4$ & $1.4084949212793479\times10^{-6}$ & $q_4$ & $2.010001314444072\times10^{-5}$\\
$p_5$ & $2.9916307576727871\times10^{-9}$ & $q_5$ & $1.1066450226912341\times10^{-7}$\\
\bottomrule
\end{tabular}
\caption{Constantes da correção $\varphi$. Todas são positivas, então o denominador
nunca se anula para $t\ge0$. Elas vêm de um aproximante de Padé, e não de ajuste numérico (Seção~\ref{sec:pade}).}
\label{tab:c}
\end{table}

Sem a correção ($\varphi\equiv0$), a fórmula \eqref{eq:formula} não tem nenhuma
constante empírica e erra no máximo $0{,}7\%$.

\section{Dedução}

\subsection{O perímetro como harmônico toroidal}

\begin{proposicao}[identidade exata]\label{prop:toro}
Para $a\ge b>0$,
\begin{equation}\label{eq:toro}
P=2\pi\sqrt{ab}\;P_{1/2}\!\left(\frac{a^2+b^2}{2ab}\right)
=2\pi\sqrt{ab}\;P_{1/2}(\cosh\xi),\qquad \xi=\ln\frac ab,
\end{equation}
em que $P_{1/2}$ é a função de Legendre de grau $\tfrac12$, também chamada de função toroidal.
\end{proposicao}
\begin{proof}
Parametrizando a elipse por $(a\cos\theta,b\sin\theta)$,
$P=\int_0^{2\pi}\sqrt{a^2\sin^2\theta+b^2\cos^2\theta}\,d\theta
=\int_0^{2\pi}\sqrt{A-B\cos2\theta}\,d\theta$, com $A=\frac{a^2+b^2}{2}$ e $B=\frac{a^2-b^2}{2}$.
Pela periodicidade, $P=2\int_0^{\pi}\sqrt{A+B\cos\phi}\,d\phi$. Como $A^2-B^2=a^2b^2$, temos
$\sqrt{A+B\cos\phi}=\sqrt{ab}\,\big(z+\sqrt{z^2-1}\cos\phi\big)^{1/2}$ com
$z=\frac{a^2+b^2}{2ab}$. A primeira integral de Laplace,
$P_\nu(z)=\frac1\pi\int_0^\pi\big(z+\sqrt{z^2-1}\cos\phi\big)^\nu d\phi$, dá
\eqref{eq:toro}. Por fim, $\cosh(\ln\frac ab)=\frac12\big(\frac ab+\frac ba\big)=z$.
\end{proof}

\begin{observacao}
A identidade \eqref{eq:toro} já é conhecida; ela aparece, por exemplo, em P.~Abbott,
\emph{On the Perimeter of an Ellipse}, The Mathematica Journal 11:2 (2009).
O que este trabalho traz de novo é a aproximação, nas seções seguintes.
\end{observacao}

A variável $\xi=\ln(a/b)$ tem uma interpretação natural: o círculo corresponde a
$\xi=0$ e a elipse degenerada a $\xi\to\infty$. Além disso, $h=\tanh^2(\xi/2)$.

\subsection{Aproximação uniforme de Olver: as funções de Bessel}

Para funções de Legendre $P_\nu(\cosh\xi)$, Olver construiu expansões
uniformes em $\xi$ cujo termo principal é
$\big(\frac{\xi}{\sinh\xi}\big)^{1/2}I_0\big((\nu+\tfrac12)\xi\big)$
(ver DLMF, \S14.15(iv)). Com $\nu=\frac12$, temos $\nu+\frac12=1$, e isso sugere
\begin{equation}\label{eq:olver1}
P\approx2\pi\sqrt{ab}\,\sqrt{\frac{\xi}{\sinh\xi}}\;I_0(\xi).
\end{equation}
Essa fórmula é exata no círculo, mas para $b\to0$ ela tende a $2\sqrt\pi\,a\approx3{,}545\,a$
em vez de $4a$, um erro de $11\%$. Acrescentamos então um termo de correção da forma usual
dessas expansões, proporcional a $\xi I_1(\xi)$:
\begin{equation}\label{eq:olver2}
P\approx2\pi\sqrt{ab}\,\sqrt{\frac{\xi}{\sinh\xi}}\;\Big[I_0(\xi)+\xi I_1(\xi)\,B(\xi)\Big],
\qquad B(\xi)=\frac{\xi\coth\xi-1}{8\xi^2}.
\end{equation}
Esta versão de dois termos não usa nenhuma constante empírica e erra no máximo
$0{,}7\%$ em todo o intervalo.

\subsection{O coeficiente exato e a correção \texorpdfstring{$\varphi$}{phi}}

Como $I_1>0$ para $\xi>0$, existe uma única função $c^*(\xi)$ que torna a
expressão exata:
\[
P=2\pi\sqrt{ab}\,\sqrt{\tfrac{\xi}{\sinh\xi}}\,\big[I_0(\xi)+\xi I_1(\xi)\,c^*(\xi)\big].
\]
Escrevemos $c^*(\xi)=B(\xi)\,(1+\varphi)$. Numericamente, $\varphi$ é pequena em
todo o intervalo: vale $0$ em $\xi=0$, chega a um máximo de $\approx0{,}085$ perto de $\xi\approx5$ e
tende a uma constante quando $\xi\to\infty$.

\begin{proposicao}\label{prop:par}
$\varphi$ é uma função par de $\xi$. Portanto ela é função de $t=\xi^2$.
\end{proposicao}
\begin{proof}
O perímetro é simétrico em $(a,b)$, e trocar $a$ com $b$ equivale a trocar $\xi$ por $-\xi$.
Além disso, $\sqrt{ab}$, $\xi/\sinh\xi$, $I_0(\xi)$, $\xi I_1(\xi)$ e $B(\xi)$ são todos pares em $\xi$.
\end{proof}

\paragraph{Série perto do círculo.} Calculando $\varphi$ com 320 dígitos e
derivando numericamente, obtém-se
\begin{equation}
\varphi(t)=\frac{7}{320}\,t-0.0028416418\,t^2+0.0002738100\,t^3-2.2526429\times10^{-5}\,t^4+\cdots
\end{equation}
O primeiro coeficiente coincide com $7/320$ em todos os dígitos calculados.

\paragraph{Limite no extremo degenerado.} Quando $b\to0$, temos $P\to4a$ com
erro exponencialmente pequeno em $\xi$. Além disso, $\sqrt{\xi/\sinh\xi}\approx\sqrt{2\xi}\,e^{-\xi/2}$,
$B\approx(\xi-1)/(8\xi^2)$ e, pelas expansões assintóticas
$I_0(\xi)\sim\frac{e^\xi}{\sqrt{2\pi\xi}}\big(1+\frac1{8\xi}+\cdots\big)$ e
$I_1(\xi)\sim\frac{e^\xi}{\sqrt{2\pi\xi}}\big(1-\frac3{8\xi}+\cdots\big)$, obtém-se
\begin{equation}\label{eq:lim}
\varphi\;\longrightarrow\;L+\frac{11\kappa-1}{\xi}+O\!\left(\frac1{\xi^2}\right),
\qquad L=8\kappa-1,\qquad \kappa=\frac{2}{\sqrt\pi}-1,
\end{equation}
ou seja, $L=0.027033336764\ldots$. A constante $2/\sqrt\pi$ vem de
$\frac{P}{2\pi\sqrt{ab}\sqrt{\xi/\sinh\xi}}\sim\frac{2}{\sqrt\pi}\cdot\frac{e^\xi}{\sqrt{2\pi\xi}}$.

\subsection{Aproximante de Padé com limite exato}\label{sec:pade}

Aproximamos $\varphi(t)$ por $N(t)/D(t)$, com $N(t)=\sum_{k=1}^5p_kt^k$ e
$D(t)=1+\sum_{k=1}^5q_kt^k$. As dez incógnitas são definidas por dez equações lineares:
\begin{enumerate}
\item a série de $N-D\varphi$ se anula nas ordens $t^1,\dots,t^9$ (nove equações);
\item $p_5=L\,q_5$, para que $\varphi(\infty)=L$, o valor exato \eqref{eq:lim}.
\end{enumerate}
A solução é a Tabela~\ref{tab:c}. Em particular, $p_1=7/320$. Todos os
coeficientes saem positivos, logo $D(t)>0$ para todo $t\ge0$ (as raízes de
$D$ são $t\approx-112$, $-14{,}7$, $-11{,}3$ e $-21{,}8\pm3{,}1i$).

\subsection{Propriedades}

\begin{teorema}
A aproximação \eqref{eq:formula}:
\begin{enumerate}
\item é exata no círculo e simétrica em $(a,b)$, se escrita com $\xi=|\ln(a/b)|$;
\item tende ao valor exato $4a$ quando $b\to0$. A convergência é lenta, com erro relativo
  da ordem de $1/\ln(a/b)$, porque o termo $1/\xi$ de \eqref{eq:lim} não é imposto;
\item reproduz a série de $\varphi$ até $t^9=\xi^{18}$, o que a torna extremamente precisa
  para elipses moderadas.
\end{enumerate}
\end{teorema}
\begin{proof}
(1) Para $\xi\to0$, todos os fatores tendem a $1$, exceto $2\pi\sqrt{ab}\to2\pi a$; a
simetria segue da Proposição~\ref{prop:par}.
(2) $\varphi\to L$ e, por construção de $L$, o termo dominante da assintótica coincide
com o de $4a$.
(3) Vale por construção.
\end{proof}

\section{Verificação numérica}

Comparamos com $P=4aE(e)$ calculado com 40 dígitos (\texttt{mpmath}). A identidade
\eqref{eq:toro} também foi verificada numericamente.

\begin{table}[H]
\centering
\begin{tabular}{cccc}
\toprule
$b/a$ & nova \eqref{eq:formula} & Olver 2 termos \eqref{eq:olver2} & Ramanujan II\\
\midrule
0.9  & $4.4\times10^{-27}$ & $5.6\times10^{-8}$ & $3.8\times10^{-18}$\\
0.7  & $6.2\times10^{-24}$ & $7.1\times10^{-6}$ & $7.0\times10^{-13}$\\
0.5  & $1.4\times10^{-21}$ & $9.0\times10^{-5}$ & $4.6\times10^{-10}$\\
0.3  & $1.7\times10^{-16}$ & $6.1\times10^{-4}$ & $7.4\times10^{-8}$\\
0.2  & $4.9\times10^{-14}$ & $1.4\times10^{-3}$ & $8.5\times10^{-7}$\\
0.1  & $3.0\times10^{-11}$ & $3.3\times10^{-3}$ & $1.2\times10^{-5}$\\
0.05 & $1.9\times10^{-9}$  & $5.0\times10^{-3}$ & $5.2\times10^{-5}$\\
0.01 & $4.2\times10^{-7}$  & $6.8\times10^{-3}$ & $2.4\times10^{-4}$\\
0.001& $1.4\times10^{-5}$  & $6.8\times10^{-3}$ & $3.8\times10^{-4}$\\
$10^{-6}$ & $2.8\times10^{-4}$ & $5.5\times10^{-3}$ & $4.0\times10^{-4}$\\
\bottomrule
\end{tabular}
\caption{Erro relativo (em módulo) por razão de semieixos.}
\end{table}

\begin{table}[H]
\centering
\begin{tabular}{lccc}
\toprule
faixa de $b/a$ & nova & Olver 2 termos & Ramanujan II\\
\midrule
$[0.2,\,1]$       & $4.7\times10^{-14}$ & $1.4\times10^{-3}$ & $8.4\times10^{-7}$\\
$[0.05,\,0.2]$    & $1.8\times10^{-9}$  & $5.0\times10^{-3}$ & $5.2\times10^{-5}$\\
$[0.01,\,0.05]$   & $3.8\times10^{-7}$  & $6.8\times10^{-3}$ & $2.3\times10^{-4}$\\
$[10^{-3},\,10^{-2}]$ & $9.1\times10^{-6}$ & $7.0\times10^{-3}$ & $3.7\times10^{-4}$\\
$[10^{-6},\,10^{-3}]$ & $2.8\times10^{-4}$ & $6.8\times10^{-3}$ & $4.0\times10^{-4}$\\
$<10^{-6}$        & $6.1\times10^{-4}$ & $5.5\times10^{-3}$ & $4.0\times10^{-4}$\\
\bottomrule
\end{tabular}
\caption{Erro relativo máximo por faixa. A nova fórmula só perde para Ramanujan II
em elipses extremamente achatadas ($b/a<10^{-6}$).}
\end{table}

\begin{figure}[H]
\centering
\begin{tikzpicture}
\begin{axis}[
  width=0.85\textwidth, height=7cm,
  xlabel={$b/a$}, ylabel={erro relativo $|\,\cdot\,|$},
  ymode=log, ymin=1e-30, ymax=1e-1, xmin=0, xmax=1,
  legend pos=south west, grid=major]
\addplot[blue, thick] table[x=x, y=nova] {erros.dat};
\addlegendentry{nova (Eq.~\ref{eq:formula})}
\addplot[green!50!black, thick, dotted] table[x=x, y=olver] {erros.dat};
\addlegendentry{Olver 2 termos}
\addplot[red, thick, dashed] table[x=x, y=ram2] {erros.dat};
\addlegendentry{Ramanujan II}
\end{axis}
\end{tikzpicture}
\caption{Erro relativo em função de $b/a$.}
\end{figure}

\section{Observações}

\begin{itemize}
\item \textbf{Área matemática.} A dedução usa ferramentas da física-matemática: harmônicos toroidais
(funções de Legendre de grau semi-inteiro, que aparecem em problemas de potencial
com simetria de toro), as integrais de Laplace e a teoria de aproximações uniformes de Olver,
que leva às funções de Bessel modificadas.
\item \textbf{Uso prático.} $I_0$ e $I_1$ estão em quase todas as bibliotecas
científicas (\texttt{scipy.special.i0}, \texttt{i1}; Excel: \texttt{BESSELI}), o que
torna a fórmula fácil de implementar.
\item \textbf{Limitação.} Para $b/a<10^{-6}$, o erro (até $6\times10^{-4}$) supera o de
Ramanujan II e decai só como $1/\ln(a/b)$.
\item \textbf{Ineditismo.} A identidade \eqref{eq:toro} é conhecida. Em buscas na
literatura (novembro de 2026), não encontramos aproximações do perímetro baseadas
nas funções de Bessel $I_0$ e $I_1$ ou na aproximação de Olver para $P_{1/2}$. A busca não foi
exaustiva.
\end{itemize}

\end{document}
